% ============================================================================
% Hypothesis H1: climate robustness of the site ranking
% Pre-registered decision rule (OQ-050, H1). Fixed before any real-country
% result was produced. Notation follows docs/METHODOLOGY.md (M-F7-01 to M-F7-07).
% Citation keys marked "verify" must be checked before submission.
% ============================================================================

\section{Hypothesis H1: Climate Robustness of the Site Ranking}
\label{sec:h1}

\subsection{Purpose and Position in the Research Design}
\label{subsec:h1_purpose}

Hypothesis H1 asks whether the spatial ranking of candidate sites produced by
GeoFREA holds across climate futures. It examines two distinct properties.
The first is \textit{core stability}: whether the sites that are robust under
one emission scenario are largely the same sites that are robust under the
others. The second is \textit{climate vulnerability}: whether a relevant share
of the sites that are optimal under present conditions lose their position once
climate change alone is introduced.

The two properties answer different questions and are therefore tested with
independent statistics. Core stability compares the Shared Socioeconomic
Pathways (SSPs) with one another; climate vulnerability compares the present
with each SSP. H1 is restricted to climate uncertainty: the comparison between
the nominal ranking and the ranking under combined climate and
techno-economic uncertainty is the object of H3
(Section~\ref{sec:h3}).

The design follows the Decision Making under Deep Uncertainty (DMDU)
perspective, in which futures are explored without assigning them probabilities
\cite{lempert2003,marchau2019}. Emission scenarios and climate models are
treated as plausible, unweighted futures, while techno-economic parameters are
treated as a sampled probabilistic space. Robustness metrics differ in what they
capture and may disagree \cite{herman2015,kwakkel2016,mcphail2018}; for this
reason one primary statistic is fixed for each sub-hypothesis, and every other
metric is reported as descriptive evidence without a verdict.

\subsection{Notation}
\label{subsec:h1_notation}

Let $\mathcal{I}$ be the set of candidate cells that enter the robustness
analysis for one country and one technology (the F7 set): cells of the central
land scenario that are feasible in the nominal future and have a valid climate
forcing in every member of the core window. Cells classified as
\textit{climate-fragile} or \textit{climate-data-invalid} are reported
separately and do not enter $\mathcal{I}$.
Let $\mathcal{I}_0$ be $\mathcal{I}$ together with the climate-fragile cells: the
cells feasible in the nominal future. The nominal set $S^{\mathrm{nom}}_{k}$ is
drawn from $\mathcal{I}_0$, so that a present-optimal site that becomes
infeasible under climate change is retained in H1b.

Climate members are $m=(g,\sigma,w)$, with $g$ one of six CMIP6 general
circulation models, $\sigma\in\Sigma=\{\mathrm{SSP1\text{-}2.6},
\mathrm{SSP3\text{-}7.0},\mathrm{SSP5\text{-}8.5}\}$ and $w$ the time window;
the core window is 2041--2070. The reference climate is $m_0$. Techno-economic
parameter vectors are $s=1,\dots,S$, drawn by Latin hypercube sampling, and
$s_0$ is the nominal vector. A future is $f=(m,s)$, and the nominal present
future is $f_0=(m_0,s_0)$.

For cell $i$ in future $f$, the relative regret is
\begin{equation}
    r_{i,f}=\frac{\mathrm{LCOE}_{i,f}-L^{*}_{f}}{L^{*}_{f}},
    \label{eq:h1_relative_regret}
\end{equation}
where $L^{*}_{f}$ is the reference LCOE of future $f$ over $\mathcal{I}$
(the minimum, or the low quantile $q_{\mathrm{ref}}$ fixed in OQ-020). The
relative form removes cost shifts common to every cell, so that the regret
reflects differences between sites rather than the general level of costs.

Two scenario-level scores are defined for each SSP $\sigma$.

\paragraph{Robust score per SSP.}
\begin{equation}
    \mathrm{MR}_{i,\sigma}
    =\max_{g}\;
    \mathrm{P}_{90,\,s\geq 1}\!\left(r_{i,(g,\sigma,w_{\mathrm{core}}),s}\right),
    \label{eq:h1_mr_ssp}
\end{equation}
the worst case over climate models of the 90th percentile over parameter
samples. It has the same structure as the primary robustness metric of the
thesis, restricted to one SSP.

\paragraph{Climate-only score per SSP.}
\begin{equation}
    R^{\mathrm{clim}}_{i,\sigma}
    =\max_{g}\; r_{i,\left((g,\sigma,w_{\mathrm{core}}),\,s_0\right)},
    \label{eq:h1_clim_ssp}
\end{equation}
the worst case over climate models with all techno-economic parameters held at
their nominal values. It isolates the effect of climate.

\paragraph{Ranks and top-$k$ sets.}
Ranks are assigned in ascending order of score (rank 1 is best), with the
deterministic tie convention recorded in the F7 phase record, applied
identically in every scenario. Let $N=|\mathcal{I}|$ and
\begin{equation}
    k=\left\lceil p_k\,N\right\rceil,
    \label{eq:h1_k}
\end{equation}
with $p_k$ the top-$k$ share fixed in OQ-021. Because $k$ scales with the
number of candidates, the same rule applies to countries whose candidate sets
differ by two orders of magnitude. The sets used by H1 are:
\begin{itemize}
    \item $S^{\mathrm{nom}}_{k}$: the $k$ cells of $\mathcal{I}_0$ with the lowest
    $\mathrm{LCOE}_{i,f_0}$ (nominal present ranking);
    \item $S^{\sigma}_{k}$: the $k$ cells with the lowest
    $\mathrm{MR}_{i,\sigma}$;
    \item $r^{\mathrm{clim}}_{i,\sigma}$: the rank of cell $i$ by
    $R^{\mathrm{clim}}_{i,\sigma}$.
\end{itemize}

\subsection{H1a: Core Stability Across Emission Scenarios}
\label{subsec:h1a}

\paragraph{Statement.}
A majority of the sites that are robust under one emission scenario are robust
under all three.

\paragraph{Primary statistic.}
The core share is the fraction of the top-$k$ set common to the three SSPs:
\begin{equation}
    C_k=\frac{1}{k}\left|\,\bigcap_{\sigma\in\Sigma} S^{\sigma}_{k}\right|.
    \label{eq:h1a_core}
\end{equation}

\paragraph{Decision rule.}
\begin{equation}
    \text{H1a holds}\iff C_k\geq\alpha,\qquad \alpha=0.60.
    \label{eq:h1a_rule}
\end{equation}

\paragraph{Reference value.}
If the three rankings were independent, the expected size of the intersection
of three random $k$-subsets of $N$ cells would be $k\,(k/N)^{2}$, so that
$\mathbb{E}[C_k]\approx p_k^{2}$. For $p_k=0.10$ this is $0.01$; the threshold
$\alpha=0.60$ requires a majority of shared sites, sixty times the chance
level. The reference value is reported with every result.

\subsection{H1b: Climate Vulnerability of Present-Optimal Sites}
\label{subsec:h1b}

\paragraph{Statement.}
At least one in five of the sites that are optimal under present conditions
lose their position under at least one emission scenario when only the climate
changes.

\paragraph{Vulnerability indicator.}
A nominal top-$k$ site is vulnerable if it falls outside the top $k+m$ of the
climate-only ranking in at least one SSP:
\begin{equation}
    I^{\mathrm{vuln}}_{i}
    =\mathbb{1}\!\left[\exists\,\sigma\in\Sigma:\;
    r^{\mathrm{clim}}_{i,\sigma}>k+m\right],
    \qquad i\in S^{\mathrm{nom}}_{k},
    \qquad m=k.
    \label{eq:h1b_indicator}
\end{equation}
A climate-fragile site in $S^{\mathrm{nom}}_{k}$ has no climate-only rank and is
counted as vulnerable.
The tolerance $m=k$ means that a site counts as vulnerable only when it leaves
the top $2k$. Movements across the top-$k$ boundary that stay within $2k$ are
treated as boundary fluctuation, because cells near the cut-off differ by
small LCOE margins and would otherwise be classified as vulnerable by
construction. The tolerance scales with $k$ for the same reason as $k$ scales
with $N$.

\paragraph{Primary statistic and decision rule.}
\begin{equation}
    V_k=\frac{1}{k}\sum_{i\in S^{\mathrm{nom}}_{k}} I^{\mathrm{vuln}}_{i},
    \qquad
    \text{H1b holds}\iff V_k\geq\beta,\qquad \beta=0.20.
    \label{eq:h1b_rule}
\end{equation}

\paragraph{Magnitude.}
Prevalence and magnitude are reported separately. For every vulnerable site
the analysis reports the worst rank displacement
$\Delta r_i=\max_{\sigma}\left(r^{\mathrm{clim}}_{i,\sigma}-r_{i,f_0}\right)$,
the SSP and climate model responsible, and $\max_{\sigma}R^{\mathrm{clim}}_{i,\sigma}$.
No threshold is applied to the magnitude, so that the verdict does not depend
on a second, unanchored parameter.

\paragraph{Independence from H1a.}
H1a compares the SSPs with one another under combined uncertainty; H1b compares
the present with each SSP under climate uncertainty alone. Neither statistic is
a function of the other, and all four combinations of outcomes are possible.

\subsection{Verdict}
\label{subsec:h1_verdict}

The verdict is issued for each country and technology:
\begin{itemize}
    \item \textit{supported}: H1a and H1b hold;
    \item \textit{partially supported}: exactly one holds;
    \item \textit{not supported}: neither holds.
\end{itemize}
The combination of the six country-technology verdicts into the cross-country
assessment follows the rule of the synthesis table (T-R8), fixed together with
the rules of H2 to H5.

\subsection{Descriptive Evidence (No Verdict)}
\label{subsec:h1_descriptive}

The following metrics are reported to characterise the result; none of them
changes the verdict.
\begin{itemize}
    \item \textit{Pairwise Jaccard overlap} of the SSP top-$k$ sets,
    $J_k(\sigma,\sigma')=|S^{\sigma}_{k}\cap S^{\sigma'}_{k}|\,/\,
    |S^{\sigma}_{k}\cup S^{\sigma'}_{k}|$, which shows which pair of scenarios
    drives a low core share \cite{fagin2003}.
    \item \textit{Rank-biased overlap} (RBO) between the SSP rankings
    \cite{webber2010}, with persistence $p=1-1/k$ so that the expected
    evaluation depth equals the top-$k$ set. With the commonly used
    $p=0.9$ the metric would weigh only the first ten or so positions, which is
    a negligible fraction of the candidate set of a large country. The finite,
    truncated estimator is used and documented.
    \item \textit{Rank displacement, median rank and interquartile range} of the
    nominal top-$k$ sites across SSPs. A small dispersion is not by itself
    evidence of robustness, since a site may be consistently poor.
\end{itemize}

\subsection{Sensitivity Analyses}
\label{subsec:h1_sensitivity}

The verdict is issued from the primary configuration only. The following
variants are reported alongside it; a change of verdict in any of them is
reported as a fragility of the result, not as a revised verdict.
\begin{enumerate}
    \item Without IPSL-CM6A-LR, the climate model kept as a recorded exception
    to the transient climate response screening.
    \item The 2071--2100 window, in which the SSPs diverge more strongly.
    \item The 0.1-degree cell scale (modifiable areal unit problem).
    \item Top-$k$ share $p_k\in\{0.05,0.10,0.20\}$ and tolerance
    $m\in\{k/2,\,k,\,2k\}$, reporting $C_k$ and $V_k$ in each case.
\end{enumerate}

\subsection{Worked Example}
\label{subsec:h1_example}

Consider a country with $N=1{,}000$ cells in $\mathcal{I}$ and $p_k=0.05$, so
$k=50$ and $m=50$.

\paragraph{H1a.}
Suppose 36 cells belong to the robust top-50 of all three SSPs. Then
$C_{50}=36/50=0.72\geq0.60$ and H1a holds. Under independence the expected
core share would be about $0.0025$.

\paragraph{H1b.}
Suppose a site ranked 5th in the nominal ranking falls to rank 78 in the
climate-only ranking of SSP5-8.5. Since $78\leq k+m=100$, it is not
vulnerable: it left the top-50 but stayed within the top-100. Another site,
ranked 12th, falls to rank 340 under SSP3-7.0; it is vulnerable. If 11 of the
50 nominal sites are vulnerable, $V_{50}=0.22\geq0.20$ and H1b holds.

\paragraph{Verdict.}
Both sub-hypotheses hold, so H1 is supported for this country and technology:
most robust sites are the same under every emission scenario, and more than
one in five of the sites that are optimal today lose their position under at
least one scenario through climate change alone.

\subsection{Parameters}
\label{subsec:h1_parameters}

\begin{table}[htbp]
    \centering
    \caption{Pre-registered parameters of H1.}
    \label{tab:h1_parameters}
    \small
    \renewcommand{\arraystretch}{1.2}
    \begin{tabularx}{\textwidth}{
        @{}
        >{\raggedright\arraybackslash}p{2.4cm}
        >{\raggedright\arraybackslash}p{2.6cm}
        >{\raggedright\arraybackslash}X
        @{}
    }
        \toprule
        \textbf{Parameter} & \textbf{Value} & \textbf{Role and justification} \\
        \midrule
        $\alpha$ & 0.60 & Minimum core share for H1a; a majority of shared
        robust sites, far above the chance level $p_k^{2}$. \\
        \addlinespace
        $\beta$ & 0.20 & Minimum share of vulnerable nominal sites for H1b;
        one in five of the present-optimal sites. \\
        \addlinespace
        $m$ & $k$ & Tolerance beyond the top-$k$ boundary; separates
        boundary fluctuation from loss of position. \\
        \addlinespace
        $k$ & $\lceil p_k N\rceil$ & Top-$k$ size scaled to the candidate set;
        $p_k$ fixed in OQ-021. \\
        \addlinespace
        $p$ (RBO) & $1-1/k$ & Descriptive metric only; expected depth equal
        to $k$. \\
        \addlinespace
        $q_{\mathrm{ref}}$ & OQ-020 & Reference quantile of $L^{*}_{f}$. \\
        \bottomrule
    \end{tabularx}
\end{table}

The thresholds $\alpha$ and $\beta$ are decision thresholds of this thesis, not
values established by the literature, which provides robustness metrics and the
recommendation to test sensitivity to them but no universal cut-offs
\cite{mcphail2018}. Their defence rests on three elements: they were fixed
before any real-country result was produced, they are interpretable (a
majority; one in five), and the sensitivity curves of
Section~\ref{subsec:h1_sensitivity} show how the conclusions depend on them.

\subsection{Limitations}
\label{subsec:h1_limitations}

\paragraph{Cut-off dependence.}
Top-$k$ membership is sensitive to small differences near the boundary. The
tolerance $m=k$ addresses this for H1b; for H1a the variation of $p_k$ and the
pairwise Jaccard curves show the extent of the effect.

\paragraph{Climate signal.}
Climate enters GeoFREA through annual-mean change factors of the mean resource,
with near-surface wind ratios applied at hub height and a fixed Weibull shape.
These simplifications may understate the spatial differentiation of climate
impacts, which would bias H1a towards stability and H1b towards rejection.

\paragraph{Ties.}
Cells with nearly identical scores can exchange positions arbitrarily. The
deterministic tie convention is applied identically in every scenario; where
ties are frequent, results are also reported with near-ties treated as
equivalent \cite{corsi2024}. % verify reference

\paragraph{Scenario interpretation.}
The SSPs are not equally probable outcomes. Results are reported per scenario
and as worst cases, without averaging across scenarios
\cite{marchau2019,kwakkel2016}.
